\usepackage[justification=centering]{caption}
% \usepackage{babel}
%\usepackage{times} \usepackage{graphicx}
% ---------------------------------------------------------------------
% LƯU Ý THỨ TỰ NẠP GÓI FONT:
%   mathptmx (font toán kiểu Times) BẮT BUỘC nạp TRƯỚC fontspec,
%   vì mathptmx có lệnh \def\rmdefault{ptm} -> nếu nạp sau fontspec thì nó sẽ
%   ghi đè \setmainfont và toàn bộ văn bản sẽ bị LaTeX thay bằng font Latin Modern
%   thay vì Times New Roman.
% ---------------------------------------------------------------------
\usepackage{mathptmx}   % same Time New Roman (nạp TRƯỚC fontspec - xem lưu ý trên)
\usepackage{fontspec}
\setmainfont{Times New Roman}
\usepackage{amsmath} 
\usepackage{indentfirst}

%\renewcommand{\rmdefault}{phv} % Arial \renewcommand{\sfdefault}{phv} % Arial \usepackage[fontsize=13pt]{scrextend}


\usepackage{fancyhdr}
% \usepackage[utf8]{inputenc}
\usepackage[vietnamese, english]{babel}
\usepackage{titlesec}

\usepackage{titletoc}
\usepackage{listings}
\usepackage[bookmarks=true]{hyperref}
\usepackage[left=3cm,right=2cm,top=2.5cm,bottom=3cm]{geometry}
\usepackage{graphicx}
\usepackage{hyperref}
\usepackage{tikz}
\usepackage{varwidth}
\usepackage{graphicx}
\usepackage{float}
\usepackage{color}
\usepackage{multirow}
\usepackage{booktabs}
% Gói thuật toán: dùng combo algorithm + algpseudocode (khớp với nội dung các chương).
% KHÔNG dùng algorithm2e vì nó xung đột với algpseudocode (\If, \Else, \EndIf).
\usepackage{algorithm}
\usepackage{algpseudocode}   % hỗ trợ môi trường algorithmic (\State, \If, \Require, ...) dùng trong nội dung các chương
\floatname{algorithm}{Thuật toán}
\counterwithin{algorithm}{chapter}   % đánh số thuật toán theo chương: 2.1
\usepackage{svg}
\usepackage{tabularx}
\usepackage{nomencl}
\usepackage{scrfontsizes}
\usepackage{longtable}
\usepackage{multicol}
\usepackage{enumitem}   % hỗ trợ tùy chỉnh danh sách: \begin{itemize}[leftmargin=*, itemsep=0.3em]
\usepackage[toc,page]{appendix}
\usepackage{adjustbox}
\usepackage{setspace}
% \usepackage{algorithm} \floatname{algorithm}{Thuật toán} \usepackage{algpseudocode}

\usepackage{makecell}
\usepackage{longtable}
\usepackage{multirow}
% \usepackage{algpseudocode}
\usepackage{setting/bkthesis}
\usepackage[toc]{appendix}

%\counterwithin{figure}{chapter}
\usepackage[ defernumbers, sortcites, backend=biber, bibstyle=authoryear, sorting=none, citestyle=numeric-comp
]{biblatex}
\usepackage{xpatch}
\usepackage{tikz}
\usetikzlibrary{shapes.geometric}
\usepackage{forest}   % vẽ cây dữ liệu (Data Tree) trong chương Phương pháp đề xuất
\usetikzlibrary{arrows.meta,arrows}
\titlelabel{\thetitle. }

% \usepackage{csquotes}
% \DeclareQuoteStyle{vietnamese}
% \UseLanguage{Vietnamese}

\makeatletter
\input{numeric.bbx}
\makeatother

\xpatchbibmacro{date+extrayear}{%
  \printtext[parens]%
}{%
  \setunit{\addperiod\space}%
  \printtext%
}{}{}

% \usepackage{biblatex}
\addbibresource{refs.bib}

% \renewcommand\labelitemi{--}

\setlength{\parskip}{6pt}

\usetikzlibrary{calc}
\setlength{\parindent}{10mm}
\renewcommand{\baselinestretch}{1.3}
\graphicspath{{images/}}
\renewcommand{\listalgorithmname}{Danh sách thuật toán}
\renewcommand{\nomname}{Danh sách ký hiệu, viết tắt}
\renewcommand\appendixtocname{Phụ lục}

% (đã gỡ \SetKwRepeat - lệnh riêng của algorithm2e; nội dung các chương không dùng)

% \algdef{SE}[DOWHILE]{Do}{doWhile}{\algorithmicdo}[1]{\algorithmicwhile\ #1}% \usepackage{setting/uetthesis}

%%% The following lines add Chapter or Appendix in front of the number
\titlecontents{chapter}%
[0pt]%
{\vspace{1ex}}%
{\bfseries Chương \thecontentslabel\quad}%
{\bfseries}%
{\bfseries\hfill\contentspage}
%%% Initially, for the main part of the document, set the label to "Chapter"
\let\chapappname\chaptername{}

\setcounter{tocdepth}{4}
\setcounter{secnumdepth}{4}

% \thesislayout hyper setup
\definecolor{dkgreen}{rgb}{0,0.6,0}
\definecolor{gray}{rgb}{0.5,0.5,0.5}
\definecolor{mauve}{rgb}{0.58,0,0.82}

% setup code area as listings
\lstset{frame=tb, language=Java, xleftmargin=10pt,xrightmargin=10pt,
%   framesep=2cm, aboveskip=2cm, belowskip=2cm,
  showstringspaces=false, columns=flexible, basicstyle={\small\ttfamily}, numbers=left, numberstyle=\tiny\color{gray},
  keywordstyle=\color{blue}, commentstyle=\color{dkgreen}, stringstyle=\color{mauve}, breaklines=true,
  breakatwhitespace=true, tabsize=2 }

\newcommand{\tcode}[1]{{\ttfamily\selectfont#1}}

\colorlet{punct}{red!60!black}
\definecolor{delim}{RGB}{20,105,176}
\colorlet{numb}{magenta!60!black}




\renewcommand\lstlistlistingname{Danh sách đoạn mã}
\renewcommand{\lstlistingname}{Đoạn mã}
\renewcommand\thechapter{\arabic{chapter}}
\renewcommand\thesection{\thechapter.\arabic{section}} 
\renewcommand\thesubsection{\thesection.\arabic{subsection}}
\renewcommand\thesubsubsection{\thesubsection.\arabic{subsubsection}}
\renewcommand{\thetable}{\thechapter.\arabic{table}} 
\renewcommand{\thefigure}{\thechapter.\arabic{figure}}

\newcolumntype{R}{>{\raggedleft\arraybackslash}X}
% L nhận tham số độ rộng cột, dùng như L{4cm} trong các bảng của nội dung chương
\newcolumntype{L}[1]{>{\raggedright\arraybackslash}p{#1}}
\newcolumntype{F}[1]{>{\raggedleft\let\newline\\\arraybackslash\hspace{0pt}}m{#1}}

\newtheorem{definition}{\indent Định nghĩa}[chapter]
